\section{Trabalhos Futuros}
\label{sec:trabalhos-futuros}

A evolução mais diretamente relacionada à modelagem é o vínculo entre o snapshot e a solicitação ou a revisão que o originou (Seção~\ref{sec:snapshot-auditoria}). A inclusão dessa chave estrangeira em SpecialistSnapshot permitiria que a reversão fosse aplicada a uma alteração específica no meio do histórico, e não apenas à última alteração aprovada (Seções~\ref{sec:snapshot-auditoria} e \ref{sec:reversao}), tornando explícita a associação hoje existente apenas por proximidade temporal e autoral. Na mesma direção, uma política de retenção e de exclusão lógica preservaria o histórico mesmo diante da remoção física de solicitações, hoje apagado pelo cascade="all, delete-orphan" (Seção~\ref{sec:snapshot-auditoria}).

No processo de integração contínua, a ativação das verificações de black, mypy e ruff como bloqueantes é uma evolução natural, condicionada à padronização prévia do código legado (Seção~\ref{sec:controle-versao-integracao-continua}). Enquanto o job de lint permanecer informativo, a garantia de qualidade recai exclusivamente sobre os testes.

No escopo funcional, duas extensões se destacam: a adição de notificações — por e-mail ou in-app — para solicitante e avaliador nos eventos do fluxo, e o suporte a múltiplas solicitações simultâneas para um mesmo especialista, com regras explícitas de conflito, uma vez que o módulo atual restringe o fluxo a uma solicitação em andamento por especialista (Seção~\ref{sec:modulo-solicitacoes}). Uma terceira extensão possível é a aplicação do mecanismo de snapshot e de linha do tempo aos demais tipos de solicitação — pontuação e recurso —, caso o domínio venha a exigir rastreabilidade equivalente à da mudança de categoria (Seção~\ref{sec:snapshot-auditoria}).

Por fim, uma evolução de natureza arquitetural seria a avaliação do Event Sourcing como referência para um histórico imutável mais robusto (Seção~\ref{subsec:snapshot-reversao}). Sua adoção efetiva exigiria uma reestruturação profunda do modelo de persistência, incompatível com a estratégia de Additive Change adotada (Seção~\ref{sec:additive-change}), mas o conceito permanece como direção possível para evoluções que exijam um histórico integralmente derivado de eventos.